\section{Motivation}

Fluorescence microscopy is one of the main tools used to study the structure and dynamics of living biological systems because it allows specific cellular and tissue components to be visualised with high contrast~\cite{...}. In developmental and cardiac imaging, this is particularly valuable since many biologically relevant structures are small, densely arranged, and embedded in complex three-dimensional environments~\cite{...}. However, the quality of fluorescence microscopy data is often limited by optical blur, anisotropic resolution, noise, and geometric distortion, all of which reduce the visibility of fine anatomical detail~\cite{...}.

These limitations are especially important in volumetric imaging of zebrafish cardiac tissue. The zebrafish is a widely used model organism for studying vertebrate heart development because of its optical transparency, rapid development, and suitability for live imaging~\cite{...}. At the same time, the myocardium presents a difficult imaging target. Its geometry is curved and highly structured, and relevant features can be small enough that even moderate degradation affects how clearly they can be observed or measured. As a result, there is a strong need for restoration methods that improve image quality while preserving biologically meaningful morphology.

Classical deconvolution methods address this problem by modelling image formation through the point spread function (PSF) and attempting to invert its effect~\cite{...}. In practice, however, these approaches often depend on accurate microscope calibration, reliable metadata, and a well-characterised PSF. Such conditions are not always easy to guarantee, especially when imaging conditions vary across datasets or when distortions are not purely optical. This creates an opportunity for data-driven methods that can learn image restoration directly from examples, without requiring explicit PSF estimation in every case~\cite{...}.

This report is motivated by the idea that deep learning can provide a more flexible approach to microscopy restoration in this setting. Rather than treating blur correction, geometric correction, and resolution enhancement as completely separate problems, it becomes possible to learn a mapping from degraded fluorescence volumes to structurally improved representations that are more useful for downstream analysis. In the case of zebrafish cardiac imaging, this is not only a question of visual improvement, but also of preserving the anatomical information needed for biological interpretation.

\section{Project Specification}

The aim of this project is to develop a deep learning pipeline for the restoration of degraded three-dimensional fluorescence microscopy images of zebrafish cardiac tissue. More specifically, the project focuses on learning a representation that compensates for blur, anisotropy, and geometric distortion while improving the effective resolution of the acquired volumes. The intention is not simply to generate visually sharper images, but to produce outputs that remain structurally plausible and better suited to further biological analysis.

A central aspect of the project is that the proposed approach does not rely on explicit estimation of the PSF. Instead of following a conventional deconvolution framework, the method learns image-based restoration directly from training data. This makes the pipeline less dependent on microscope-specific calibration and potentially more robust to variations in acquisition conditions. In this sense, the work explores whether a data-driven model can recover useful structural information even when the degradation process is only partially characterised.

The project is situated at the intersection of fluorescence microscopy, computational imaging, and deep learning. On the imaging side, it requires an understanding of how biological structures are represented in volumetric microscopy and how this representation is altered by physical and optical limitations. On the computational side, it involves designing and evaluating a restoration framework capable of improving image quality while preserving morphological consistency. The zebrafish myocardium provides a meaningful test case because it combines biological relevance with challenging image characteristics.

More broadly, the project contributes to the ongoing shift from explicitly model-based restoration toward learned image reconstruction methods in microscopy~\cite{...}. While classical approaches remain important and well grounded in physical theory, learning-based methods offer the possibility of handling more complex or spatially varying degradations that are difficult to model analytically. This report therefore investigates a practical restoration strategy designed for a realistic biological imaging problem rather than an idealised benchmark setting.

\section{Report Structure}

The remainder of this report is organised as follows. Chapter~2 presents the background needed to understand the biological and imaging context of the project. It introduces fluorescence microscopy in cardiac imaging, outlines the relevance of zebrafish heart morphology, and reviews the role of the point spread function, morphometric degradation, and resolution limits in microscopy image formation.

Chapter~3 describes the methodology used in this work. This includes the dataset, preprocessing steps, model design, training strategy, and evaluation framework used to assess the performance of the proposed restoration pipeline. Particular attention is given to how geometric correction and resolution enhancement are handled within the learning-based approach.

Chapter~4 presents the experimental results and discusses their interpretation. The restored outputs are analysed in relation to the original degraded images, with emphasis on image quality, structural plausibility, and the preservation of biologically relevant features. Where appropriate, the strengths and limitations of the method are discussed in relation to both classical restoration techniques and the specific challenges of zebrafish cardiac imaging.

Finally, Chapter~5 concludes the report by summarising the main findings, identifying the main limitations of the current approach, and outlining possible directions for future work.